\documentclass[10pt]{article}

\usepackage[utf8]{inputenc}

\usepackage[T1]{fontenc}

\usepackage{amsmath}

\usepackage{amsfonts}

\usepackage{amssymb}

\usepackage[version=4]{mhchem}

\usepackage{stmaryrd}

\usepackage{bbold}



\begin{document}

П. ф.- аналог понятия «свободная универсальная алгебра с одним образующим». Для любого функтора $G: \Re \rightarrow \varsigma$ и П. ф. $F$ множество естественных преобразований $\operatorname{Nat}(F, G)$ изоморфно множеству $G(A)$, где $A$ - представляющий объект. Это показывает, что П. ф. являются свободными объектами категории функторов.



Для аддитивных категорий вместо функторов со значениями в $\Im$ рассматриваются аддитивные функторы со значениями в категории абелевых групп; поэтому под П. ф. понимается аддитивный функтор, изоморфный основному аддитивному функтору.



Понятие П ф. первоначально возникло в алгебраич. геометрии (см [2]). Наиболее важными примерами П. Ф. здесь являются функторы Пикара Ріс $X / S$ и Гильберта $\mathrm{Hilb} X / S$, представимые в категории алгебраических пространств (см. [1]). Пусть $K$ - поле частных регулярного дискретного нормированного кольца $O$ с совершенным полем вычетов. Если $X_{0}$ - гладкая геометрически неприводимая собственная кривая. рода $g>0$ над $K$, то ее минимальная модель представляет функтор $Y \mapsto \operatorname{Isom}_{K}\left(Y \bigotimes_{o} K, X_{0}\right)$ из категории регулярных $O$-схем. Если $A$ - абелево многообразие над $K$, то его минимальная Нерона модель является гладкой групповой схемой $X \rightarrow \mathrm{Spec} 0$, представляющей функтор $Y \mapsto \operatorname{Hom}_{K}\left(Y \otimes_{o} K, A\right)$ из категории гладких O-схем.



Лит.: [1] А р т и н М., "Математика", 1970, т. 14, № 4, с. 3-39; [2] Грот е ндин А., Д ьё дон н е Ж., "Успехи матем. наук», 1972, т. 27 , в. 2, с. $135-48$.



ПРЕДСТАВЛЕНИЕ АЛГЕБРЫ ЛИ В В Р К Т о рн о м п р о с т p а н с т в е $V$ - гомоморфизм $\rho$ алгебры Ли $L$ над полем $k$ в алгебру Ли $\mathfrak{l l}(V)$ всех линейных преобразований пространства $V$ над $k$. Два представления $\rho_{1}: L \rightarrow g \mathfrak{l}\left(V_{1}\right) \quad$ и $\quad \rho_{2}: L \rightarrow g \mathfrak{l}\left(V_{2}{ }_{2}\right)$ наз. ә к в и в ал е н т ны м и (или и зом о р ф ным и), если существует изоморфизм $\alpha: V_{1} \rightarrow V_{2}$, для к-рого



\[

\alpha\left(\rho_{1}(l) v_{1}\right)=\rho_{2}(l) \alpha\left(v_{1}\right)

\]



при любых $l \in L, \quad v_{1} \in V_{1}$. Представление $\rho$ в $V$ наз. ко н е ч но м е р ным, если $\operatorname{dim} V<\infty$, и не п р ив о д у м ы м, если в $V$ не существ ует отличных от нуля и всего пространства подпространств, инвариантных относительно всех операторов $\rho(l), l \in L$.



При заданных представлениях $\rho_{1}: L \rightarrow \mathfrak{g l}\left(V_{1}\right)$ и $\rho_{2}: L \rightarrow \mathfrak{g l}\left(V_{2}\right)$ можно построить представления $\rho_{1} \bigoplus \rho_{2}$ (прямая сумма) и $\rho_{1} \otimes \rho_{2}$ (тензорное произведение) алгебры $L$ в пространствах $V_{1} \oplus V_{2}$ и $V_{1} \otimes V_{2}$, полагая $\left(\rho_{1} \oplus \rho_{2}\right)(l)\left(v_{1}, v_{2}\right)=\left(\rho_{1}(l) v_{1}, \rho_{2}(l) v_{2}\right)$, $\left(\rho_{1} \otimes \rho_{2}\right)(l) v_{1} \otimes v_{2}=\rho_{1}(l) v_{1} \otimes v_{2}+v_{1} \otimes \rho_{2}(l) v_{2}$ для $v_{1} \in V_{1}, v_{2} \in V_{2}, \quad l \in L$. Если $\rho$ - представление алгебры Ли $L$ в пространстве $V$, то формула



\[

<\rho^{*}(l) u, v>=-<u, \rho(l) v>

\]



определяет представление $\rho^{*}$ алгебры $L$ на сопряженном к $V$ пространстве, наз. к о н т р а г р а д и ен т ны м по отношению к $\rho$.



Каждое П. а. Ли $L$ однозначно продолжается до представления универсальной обертывающей алгебры $U(L)$; тем самым устанавливается изоморфизм категории П. а. Ли $L$ и категории модулей над $U(L)$. В частности, представлению $\rho$ алгебры $L$ соответствует идеал $\operatorname{ker} \tilde{\rho}$ в $U(L)$ - ядро его продолжения $\tilde{\rho}$. Если представление $\rho$ неприводимо, то идеал ker $\tilde{\rho}$ примитивен. Обратно, всякий примитивный идеал в $U(L)$ строится таким способом по нек-рому (вообще говоря, не единственному) неприводимому представлению $\rho$ алгебры $L$. Изучение пространства примитивных идеалов $\operatorname{Prim} U(L)$, снабженного топологией Джекобсона, является существенной частью теории П. а. Ли. Оно проведено пол- ностью в случае, когда $L$ - конечномерная разрешимая алгебра Ли, а $k$-алгебраически замкнутое поле характеристики нуль (см. [2]).



Наиболее полно изучены конечномерные представления конечномерных алгебр Ли над алгебраически замкнутым полем характеристики нуль. В случае полей $\mathbb{C}$ и $\mathbb{R}$ эти представления находятся во взаимно однозначном соответствии с аналитическими конечномерными представлениями соответствующих односвязных (комплексных или вещественных) групп Ли. В этой ситуации любое представление разрешимой алгебры Ли содержит одномерное инвариантное подпространство (см. Ли теорема). Любое представление полупростой алгебры Ли вполне приводимо, т. е. изоморфно прямой сумме неприводимых представлений. Неприводимые представления полупростой алгебры Ли $L$ полностью классифици рованы: классы изоморфных представлений взаимно однозначно соответствуют д о м ч н а н тн ы м в е с а м, т. е. весам с неотрицательными числовыми отметками, из сопряженного пространства $H^{*}$ к подалгебре Картана $H$ алгебры $L$ (см. Картана теорема о старшем векторе). Об описании строения неприводимого представления по соответствующему өму доминантному весу (его старшему весу) см. в статьях К ность веса, Характеров бормула.



Произвольный (не являющийся, вообще говоря, доминантным весом) элемент $\lambda \in H^{*}$ также определяет нек-рое неприводимое линейное представление полу простой алгебры Ли $L$ со старшим весом $\lambda$, являющееся, однако, бесконечномерным (см. Представление со старшим вектором). Соответствующие $U(L)$-модули наз. м о д у л я м и В е р м а (см. [2]). Полной классифйкации неприводимых бесконечномерных представлений полупростых алгебр Ли пока (1983) не получено.



Если $k$ - алгебраически замкнутое поле характеристики $p>0$, то неприводимые представления конечномерной алгебры Ли $L$ всегда конечномерны и их размерность ограничена константой, зависящей от $n=$ $=\operatorname{dim} L$. Если алгебра $L$ имеет $p$-структуру, то эта константа есть $p^{(n-r) / 2}$, где $r$ - минимальная возможная размерность аннулятора линейной формы на $L$ в коприсоединенном представлении [4]. Для описания множества неприводимых представлений в этом случае применяется следующая конструкция. Пусть $Z(L)-$ центр алгебры $U(L)$ и $M_{L}$ - аффинное алгебраич. многообразие (размерность $\operatorname{dim} M_{L}=n$ ), алгебра регулярных функций на к-ром совпадает с $Z(L)$ (м ного о б р а и е Ца с с е т х у з а). Отображение $\rho \mapsto \operatorname{ker}(\tilde{\rho} / Z(L))$ позволяет сопоставить каждому неприводимому представлению точку многообразия Цассенхауза. Получаемое отображение сюръективно, прообраз любой точки из $M_{L}$ конечен, а для точек открытого всюду плотного подмножества этот прообраз состоит из одного элемента [7]. Полное описание всех неприводимых представлений имеется для нильпотентных алгебр Ли (см. [8]) и нек-рых отдельных примеров [cм. [9], [10]). Получены также разнообразные результаты относительно специальных типов представлениӥ. Лum.: [1] Б у $\mathbf{p}$ б а к и Н., Групшы и алгебры Ли, пер. с франц., М., 1976-78; [2] Д и к с м ь е Ж., Универсальные

обертывающие алгебры, пер. с франц., М., 1978; [3] Д ж е к о бобертывающие алгебры, пер. с франц., М., 1964 ; [ [4] М и л в н е р A. А., "Функциональный анализ и его приложения», 1980, т. 14.

$N_{2}$ 2, с. 67-68; [5] С е р р Ж. - П., Алгебры Ли и группы Ли, $N_{2}$ 2, с. 67-68; [5] C е p p ЖЖ. -П., Алгебры Ли п группы Ли,

пер. с англ. и франц., М., 1969; [6] Теория алгебр Ли. Топология групп Ли, nep. c франц., М., 1962; [7] Z a s s e n h a u s H., "Proc. Glasgow Math. Assoc.", 1954, v. 2, p. 1-36; [8] B е й cего приложения", "1971, т. 5. Г.'2, "Функциональный анализ и J. C., "Math. Z.", 1974, Bd 140, H. 1, S. 127-49; [10] P y д aк о в А. Н., “Изв. АН СССР. Сер. матем.", 1970, т. 34, № 4,

с $735-43$. IІРЕДСТАВЛЕНИЕ АССОЦИАТИВНОЙ $\begin{aligned} & \text { А. РУдаков. } \\ & \text { АЈГЕБ- }\end{aligned}$ РЫ р а з м е р н о с т и $n$ - гомоморфизм алгебры $A$ над полем $F$ в алгебру матриц $M_{n}(F)$, т. е. сопостав-





\end{document}